\begin{pastexam}[label={px:gw3-111-ds-03}]{關務三等資訊處理 111 年 第三題}
（申論；20 分）

7 個工作的利潤及處理的最後期限如下表所示，並假設每個工作的處理時間均為一個時間單位，請求出最適排程及最大
利潤（maximal profit）為何？（20 分）
\begin{center}
\begin{tabular}{|l|*{7}{c|}}
\hline
工作編號 & 1 & 2 & 3 & 4 & 5 & 6 & 7 \\ \hline
利潤 & 40 & 15 & 60 & 20 & 45 & 10 & 55 \\ \hline
最後期限 & 2 & 4 & 3 & 2 & 1 & 3 & 1 \\ \hline
\end{tabular}
\end{center}
\end{pastexam}

\begin{pastsoln}
答案：做工作 $7, 1, 3, 2$，依序放在時段 $1, 2, 3, 4$；最大利潤 $55 + 40 + 60 + 15 = \mathbf{170}$。

這是「有期限的單位時間工作排程」（資結書第~19.4.2~節）。 時間切成時段 $1, 2, 3, \ldots$，每個時段做一件工作；
工作 $i$ 放在時段 $t$ 能如期完成的條件是 $t \le d_i$。 要選一些工作、排進不同的時段且都如期完成，使利潤總和最大。

\textbf{貪心法}：工作依利潤由大到小考慮；對每件工作，找「不超過它期限的空時段」中最晚的一個放進去；找不到
就放棄這件工作。 放最晚的空時段，是為了把前面的時段留給期限較早的工作。 期限最大是 $4$，只需要時段 $1$--$4$。

依利潤排序：$J_3$ 60、$J_7$ 55、$J_5$ 45、$J_1$ 40、$J_4$ 20、$J_2$ 15、$J_6$ 10。
\begin{center}
\small
\begin{tabular}{cccllc}
\toprule
工作 & 利潤 & 期限 & 處理 & 時段 1--4 & 累計利潤 \\
\midrule
$J_3$ & 60 & 3 & 時段 3 空，放入 & --, --, $J_3$, -- & 60 \\
$J_7$ & 55 & 1 & 時段 1 空，放入 & $J_7$, --, $J_3$, -- & 115 \\
$J_5$ & 45 & 1 & 時段 1 已占，沒有空時段，放棄 & $J_7$, --, $J_3$, -- & 115 \\
$J_1$ & 40 & 2 & 時段 2 空，放入 & $J_7$, $J_1$, $J_3$, -- & 155 \\
$J_4$ & 20 & 2 & 時段 2、1 都已占，放棄 & $J_7$, $J_1$, $J_3$, -- & 155 \\
$J_2$ & 15 & 4 & 時段 4 空，放入 & $J_7$, $J_1$, $J_3$, $J_2$ & 170 \\
$J_6$ & 10 & 3 & 時段 3、2、1 都已占，放棄 & $J_7$, $J_1$, $J_3$, $J_2$ & 170 \\
\bottomrule
\end{tabular}
\end{center}
最適排程：時段 1 做 $J_7$、時段 2 做 $J_1$、時段 3 做 $J_3$、時段 4 做 $J_2$；最大利潤 $\mathbf{170}$。

\textbf{為什麼是最大}。 一組工作排得進去（每件都能如期完成），若且唯若對每個 $t$，期限 $\le t$ 的工作不超過 $t$ 件
（資結書引理~19.8）。 資結書第~19.4.2~節用這個條件證明了「依利潤由大到小、能加就加」的貪心法得到最大利潤。 具體來說：選了 $\{J_7, J_1, J_3, J_2\}$
之後，被放棄的 $J_5$、$J_4$、$J_6$ 各自加進來都會違反條件（$J_5$ 讓期限 $\le 1$ 的有 $2$ 件；$J_4$ 讓期限 $\le 2$ 的有
$3$ 件；$J_6$ 讓期限 $\le 3$ 的有 $4$ 件），而貪心法每次留下的都是當時還能加入的利潤最大的工作。

利潤 $170$ 的工作集合只有 $\{J_1, J_2, J_3, J_7\}$ 一組；而且 $J_7$ 只能在時段 1，$J_1$ 就只能在時段 2，$J_3$ 只能在
時段 3，$J_2$ 在時段 4，所以最適排程是唯一的。
\end{pastsoln}
